% Anleitung zu den beiden TI-Nspire-Lua-Skripten
%   Tragwerksskript/Biegelinieneu.lua  und  Querschnitte/Querschnitt.lua
% Uebersetzen mit pdflatex (zweimal, wegen Inhaltsverzeichnis und Verweisen).
% Bilder: \bildplatz{dateiname}{Bildunterschrift} setzt einen Platzhalterrahmen.
%         Zum Einfuegen eines echten Bildes die Zeile durch
%         \begin{figure}[htbp]\centering\includegraphics[width=0.8\linewidth]{bilder/datei}\caption{...}\end{figure}
%         ersetzen.
\documentclass[11pt,a4paper]{article}

\usepackage[utf8]{inputenc}
\usepackage[T1]{fontenc}
\usepackage[ngerman]{babel}
\usepackage[a4paper,margin=2.4cm]{geometry}
\usepackage{amsmath}
\usepackage{amssymb}
\usepackage{array}
\usepackage{booktabs}
\usepackage{longtable}
\usepackage{graphicx}
\usepackage{xcolor}
\usepackage[autostyle,german=quotes]{csquotes}
\usepackage[hidelinks]{hyperref}

\setlength{\parindent}{0pt}
\setlength{\parskip}{0.6em}

\newcommand{\taste}[1]{\textsf{\textbf{#1}}}
\newcommand{\code}[1]{\texttt{#1}}
\newcommand{\bildplatz}[2]{%
  \begin{figure}[htbp]
    \centering
    \fbox{\parbox[c][3.0cm][c]{0.78\linewidth}{\centering\itshape Platz für ein Bild: \texttt{#1}}}
    \caption{#2}
  \end{figure}}

\title{Anleitung\\[0.3em]\large Tragwerksskript und Querschnittsskript für den TI-Nspire CX II CAS}
\author{}
\date{}

\begin{document}
\maketitle
\tableofcontents
\newpage

%=====================================================================
\section{Grundlegendes}

\subsection{Überblick}

Das Programm besteht aus drei Modulen, und deckt große teile der TM1 und TM2 Klausur ab sowie auch Teile der Statik Klausur. Dafür gibt es drei Module, also Seiten auf dem Taschenrechner. Diese sind:\\

\begin{itemize}
  \item \textbf{Tragwerksskript}: ebene Stabtragwerke und Fachwerke, Schnittgrößen,
        Verformungen, Kraftgrößenverfahren, Prinzip der virtuellen Verschiebungen, Export der Biege- und Längslinien ins CAS. Dieses Modul ist relevant für \textbf{TM1}, \textbf{TM2} und \taste{Statik}.
        
  \item \textbf{Querschnittsskript}: Querschnittswerte,
        Normalspannung,Schubspannung, Schubmittelpunkt, Torsion, Kernfläche. Vor allem relevant für \textbf{TM2} für einfache Schwerpunkt aufgaben auch \textbf{TM1}.

  \item \textbf{Spannungskreis}: ebener Spannungszustand über Schnitte eingeben, mit Tabelle,
        Spannungsscheibe und Mohrschem Kreis -- falls ihr euch auch immer in die falsche Richtung dreht. Relevant für \textbf{TM2}. 
\end{itemize}

Für die Meisten Exporte und für symbolische Berechnungen wird meines Wissen nach die
\textbf{CAS}-Version des Rechners gebraucht. Da bin ich mir aber nicht sicher.



\subsection{Grundprinzip der Bedienung}

\subsubsection{Tragwerk- und Querschnittmodul}

\begin{itemize}
  \item Gezeichnet wird mit der Maus bzw. dem Touchpad, Punkte rasten auf dem Raster ein.
  \item Ein Objekt unter dem Mauszeiger (\emph{Hover}) wird hervorgehoben. \taste{Enter} öffnet
        dessen Menü.
  \item Obermenü mit Taste \taste{o}. In Menüs bewegt man sich mit den Pfeiltasten \taste{$\uparrow$}/\taste{$\downarrow$};
        \taste{$\leftarrow$}/\taste{$\rightarrow$} blättert die Seiten. \taste{Enter} öffnet eine
        Eingabe oder schaltet einen Wert um, ein zweites \taste{Enter} übernimmt die Eingabe.
  \item \taste{Esc} bricht ab und schließt Boxen, \taste{Backspace} löscht.
  \item Alle Zahlenfelder akzeptieren Ausdrücke: \code{2/3}, \code{sqrt(2)}, \code{1e8},
        Dezimalkomma (\code{1,5}) und – im Tragwerksskript – Symbole wie \code{q} oder \code{EI}.
\end{itemize}

\bildplatz{uebersicht}{Startbildschirm der beiden Skripte}

%=====================================================================
\part{Tragwerksskript}

\section{Systemeingabe}

Die eigentliche Systemberechnung läuft über das WGV (kommt in Statik), aus diesem Grund mag die Art Systeme zu zeichnen manchmal evtl. etwas unintuitiv wirken.\\
Die zwei Grundlegenden Bausteine eines Systems sind Elemente bzw. Stäbe und Knoten.
Knoten können dabei überall platziert werden, Stäbe können nur von einem Knoten zu einem anderen Knoten platziert werden.
Alle weiteren Eigenschaften des Systems wie Lagerung, Lasten und Gelenke werden entweder Über Knoten oder Stäbe eingegeben. (Ein Einfeldträger hat also zwei Knoten und ein Stab, ein 3-Gelenk-Tragwerk hat drei knoten und zwei Stäbe.) \\
Für \textbf{Statik} sollte diese Eingabe Sinn machen.\\ 
Für \textbf{TM2} könnt ihr euch merken, das für jeden Bereich für den ihr eine neue Biegelinie ansetzten würdet, ihr auch einen neuen Stab setzten müsst. Das liegt daran, dass das Programm effektiv für alle Stabe die DGL der Biegelinie löst, die FEM bzw. das WGV (in diesem Fall zwei direkt verwandte Verfahren) bieten nur einen etwas raffinierteren Lösungsweg als den den man in der \textbf{TM2} lernt.\\
Für \textbf{TM1} fällt mir leider kein besserer Hinweis ein als einfach genügend Systeme einzugeben, biss man ein Gefühl für die Eingabe bekommt. 


\bildplatz{tragwerk-system}{Systemansicht mit Lagern, Gelenken und Lasten}

\section{Knoten}

Am Knoten hängt alles, was punktuell ist: Lager, Einzellasten, Federn und das Vollgelenk. Ihr
kommt ins Knotenmenü, indem ihr mit dem Zeiger auf den Knoten geht und \taste{Enter} drückt --
anklicken müsst ihr nicht. Das Menü hat vier Seiten, zwischen denen ihr mit
\taste{$\leftarrow$}/\taste{$\rightarrow$} blättert.

\begin{longtable}{p{0.16\linewidth}p{0.36\linewidth}p{0.42\linewidth}}
  \toprule
  \textbf{Seite} & \textbf{Zeilen} & \textbf{Bedeutung} \\
  \midrule \endhead
  1 Geometrie & X-Koordinate, Y-Koordinate, Gelenkig, Lager X, Lager Y, Lager M &
    Koordinaten dürfen auch ein Ausdruck oder ein Symbol sein. \emph{Gelenkig} macht den Knoten
    zum Vollgelenk (alle Stabenden dort momentenfrei). Die drei Lager sind einzeln schaltbar:
    X und Y zusammen ergeben ein Festlager, nur eines davon ein Rollenlager, alle drei eine
    Einspannung. \\
  2 Lasten & Last $F_x$, Last $F_y$, Last $M$, Lastwinkel &
    Einzellasten am Knoten. Der Winkel dreht das Lastpaar $F_x$/$F_y$ mit. \\
  3 Federn & Feder $c_x$, $c_y$, $c_m$, Federwinkel &
    Weg- und Drehfedern; der Winkel dreht die Federachsen. Federwerte dürfen symbolisch sein,
    z.\,B.\ \code{3*EA*l}. \\
  4 Ergebnisse & $R_\xi$, $R_\eta$, $R_m$, \code{ü\_x}, \code{ü\_y}, \code{ü\_ges}, Export &
    Auflagerreaktionen in Lagerrichtung nach der Berechnung; ganz unten könnt ihr den Knoten
    einzeln ins CAS exportieren. \\
  \bottomrule
\end{longtable}

Ein \textbf{gedrehtes Lager} (schiefe Ebene) macht ihr über den Lagerwinkel: Die Lagerrichtungen
$\xi$ und $\eta$ drehen sich mit, und die Auflagerkräfte werden auch in diesen Richtungen
ausgegeben -- ihr müsst also nichts mehr von Hand zerlegen.

\section{Stäbe}

Der Stab trägt alles, was über die Länge verteilt ist: Steifigkeiten, Streckenlasten, Temperatur
und die Gelenke an seinen Enden. Auch hier: Zeiger auf den Stab, \taste{Enter}, sieben Seiten.

\begin{longtable}{p{0.16\linewidth}p{0.38\linewidth}p{0.40\linewidth}}
  \toprule
  \textbf{Seite} & \textbf{Zeilen} & \textbf{Bedeutung} \\
  \midrule \endhead
  1 Steifigkeit und Lasten & $EI$, \emph{biegesteif}, $EA$, \emph{dehnsteif}, Last $q(x)$,
    Last $n(x)$, Momentenlast $m$, Bogenradius &
    $q$ wirkt quer, $n$ längs, $m$ ist ein Streckenmoment. \emph{biegesteif}/\emph{dehnsteif}
    machen den Stab starr (Abschnitt~\ref{sec:starr}); das zugehörige Feld ist dann gesperrt und
    zeigt \emph{starr}. Ein Bogenradius $\ne 0$ macht aus dem Stab einen Kreisbogen. \\
  2 Globale Linienlasten & Global $p_X$, $p_X$ projiziert, Global $p_Y$, $p_Y$ projiziert &
    Lasten in globaler Richtung, wahlweise auf die Stabachse projiziert -- also genau die
    Unterscheidung \enquote{Schnee auf der Dachfläche} gegen \enquote{Schnee auf dem Grundriss}. \\
  3 Trapezlasten und Temperatur & Trapez $q$ Start/Ende, Trapez $n$ Start/Ende,
    $T_{oben}$, $T_{unten}$, $\alpha_T$, Höhe $h$ &
    Die Trapezanteile addieren sich zur konstanten Last aus Seite~1. Temperatur siehe
    Abschnitt~\ref{sec:temperatur}. \\
  4 Gelenke & M-Gelenk Start/Ende, N-Gelenk Start/Ende, Q-Gelenk Start/Ende &
    Ein M-Gelenk überträgt kein Moment, ein N-Gelenk keine Normalkraft, ein Q-Gelenk keine
    Querkraft. Das Gelenk sitzt am jeweiligen Stabende, nicht am Knoten. \\
  5 Schnittgrößen & $N_a$, $Q_a$, $M_a$, $N_b$, $Q_b$, $M_b$ & Stabendgrößen nach der Berechnung. \\
  6 Auflager-/Endkräfte & $h_a$, $v_a$, $M_a$, \dots & Dieselben Endkräfte, aber in globaler
    Richtung -- praktisch fürs Freischneiden. \\
  7 Verformungen und Export & $u_a$, $v_a$, $\varphi_a$, \dots, \emph{In TI speichern} &
    Stabendverformungen; Export des einzelnen Stabes ins CAS. \\
  \bottomrule
\end{longtable}

\taste{f} dreht den Stab unter dem Zeiger um: Anfang und Ende tauschen, die Anfangs- und
Endlasten wandern mit. Das braucht ihr, wenn die lokale $x$-Richtung für den Export in die
andere Richtung laufen soll.

\bildplatz{tragwerk-stabmenue}{Stabmenü, Seite 1}

\section{Berechnung, Ansichten und Diagramme}

\taste{b} rechnet. Danach stehen euch diese Ansichten zur Verfügung:

\begin{longtable}{p{0.12\linewidth}p{0.84\linewidth}}
  \toprule
  \textbf{Taste} & \textbf{Ansicht} \\
  \midrule \endhead
  \taste{s} & Systemansicht \\
  \taste{n} & Normalkraft $N$ \\
  \taste{q} & Querkraft $Q$ \\
  \taste{m} & Moment $M$ \\
  \taste{w} & Biegelinie $w$ \\
  \taste{u} & Längsverformung $u$ \\
  \taste{e} & Explosionsansicht: alle Stäbe freigeschnitten, mit \taste{e} blättert ihr die
              Seiten durch \\
  \taste{t} & Beschriftungen in der Explosionsansicht ein- und ausblenden \\
  \taste{v} & Kinematik bzw.\ Verschiebungszustand (PvV) \\
  \taste{p} & Polplan \\
  \taste{l} & globales Gleichungssystem anzeigen \\
  \bottomrule
\end{longtable}

Die \textbf{Explosionsansicht} (\taste{e}) ist für \textbf{TM1} und \textbf{Statik} die
nützlichste: Dort seht ihr jeden Stab einzeln freigeschnitten mit allen Kräften, die an ihm
angreifen -- also genau das Bild, das ihr sonst von Hand malt. Seite~1 zeigt die Zahlenwerte,
Seite~2 dieselben Kräfte mit ihren Namen ($Gx_1$, $Gn_2$, $A_h$ \dots), so wie sie auch im
LGS-Export heißen.

\textbf{Streckenlasten stehen dort als Resultierende:} Statt der Lastfläche seht ihr einen Pfeil
auf der Wirkungslinie, beschriftet mit $R = \dots$, und die Lage hängt in der Maßkette unter dem
System -- den Hebelarm könnt ihr also direkt ablesen. Das gilt für die Querlast $q$, die Längslast
$n$ und die globalen Linienlasten $p_X$, $p_Y$; ein Streckenmoment $m$ wird zu einem Moment
$M = m\,l$ in Stabmitte. Bei einer Dreieckslast sitzt der Pfeil, wo er hingehört, nämlich bei
$x_R = S/R$ vom Stabanfang aus, und bei einer Lastfunktion wie \code{x\^{}2} rechnet das
Programm $R$ und die Lage aus der Lastfunktion selbst.

Im Obermenü (\taste{o}) stellt ihr die Darstellung ein: Seite~1 die Anzeige der Normalkraft und
den Schnittgrößenmodus (\emph{Hover} zeigt den Wert unter dem Zeiger, \emph{Start\&End} die
Stabendwerte), Seite~3 das Zahlenformat (Wurzelbrüche, nur Brüche, Dezimal), die Bemaßung, die
Punkte pro Stab und die Maximalwerte, Seite~4 die Verlaufsskalierung.

\bildplatz{tragwerk-momente}{Momentenverlauf mit Maximalwerten}

\section{Steifigkeiten und starre Stäbe}\label{sec:starr}

Jeder Stab hat ein $EA$ und ein $EI$. Ab Werk sind das $EA = 10^{10}$ und $EI = 10^{8}$, also
Werte, die praktisch starr sind -- das ist für \textbf{TM1} und \textbf{Statik} genau richtig,
weil dort nur die Kräfte interessieren. Im Obermenü Seite~1 setzt ihr beide Werte für alle Stäbe
auf einmal (\emph{Alle EA setzen}, \emph{Alle EI setzen}); der Wert gilt dann auch für Stäbe, die
ihr danach noch zeichnet, und er darf symbolisch sein (\code{EA}, \code{EI}).

Für \textbf{TM2} wird es dann interessant, wenn Stäbe wirklich starr sein sollen. Dafür gibt es
im Stabmenü Seite~1 zwei Schalter:

\begin{itemize}
  \item \emph{biegesteif}: $EI \to \infty$, der Stab verbiegt sich nicht.
  \item \emph{dehnsteif}: $EA \to \infty$, der Stab ändert seine Länge nicht.
\end{itemize}

Das wirkt überall: in den Diagrammen, in den Auflagerkräften und im Export. Intern werden die
Stabenden dabei exakt gekoppelt und die Schnittgrößen aus dem Gleichgewicht bestimmt. Eine sehr
große Ersatzsteifigkeit wird bewusst \emph{nicht} verwendet, weil das neben weichen Federn zu
Zahlenmüll führt.

Zusätzlich gibt es im Obermenü Seite~4 den Schalter \emph{Standard-EA dehnstarr} (ab Werk aus).
Damit gelten alle Stäbe als dehnsteif, bei denen ihr im Stabmenü kein eigenes $EA$ eingetragen
habt. Das ist praktisch, wenn in der Aufgabe steht \enquote{alle Stäbe dehnstarr, nur Stab~3 mit
$EA$}: Schalter an, und Stab~3 bekommt sein $EA$ im Stabmenü. Ein symbolisches $EA$ bleibt immer
dehnbar.

Der Schalter \emph{Starre Stäbe (Export)} (ebenfalls Seite~4) wählt das Verfahren für den
Rand-LGS-Export: \emph{direkt} (Standard, schneller) oder \emph{Grenzwert}. Beide liefern
dieselben Ergebnisse; halten sich starre Teile gegenseitig statisch unbestimmt, wechselt das
Programm von selbst auf den Grenzwert.

\textbf{Pendelstäbe} müsst ihr als solche kennzeichnen, nämlich mit Gelenken an beiden Enden
(M-Gelenke im Stabmenü Seite~4 oder gelenkige Knoten). Ohne Gelenke ist ein Stab immer ein
Biegestab und macht mit seinem $EI$ beim Tragverhalten mit -- das ist eine der häufigsten
Fehlerquellen bei der Eingabe.

\section{Temperatur}\label{sec:temperatur}

Im Stabmenü Seite~3 stehen $T_{oben}$, $T_{unten}$, der Ausdehnungskoeffizient $\alpha_T$ und die
Querschnittshöhe $h$. Daraus macht das Programm die zwei Anteile, die ihr aus \textbf{TM2} kennt:

\begin{align*}
  \varepsilon_T &= \alpha_T \,\frac{T_{oben} + T_{unten}}{2} && \text{(gleichmäßige Erwärmung, Längenänderung)}\\
  \kappa_T &= \alpha_T \,\frac{T_{unten} - T_{oben}}{h} && \text{(Temperaturunterschied, Krümmung)}
\end{align*}

\textbf{Aufpassen bei $h = 0$:} Dann gibt es keinen Gradienten, es wirkt nur die mittlere
Erwärmung. Wenn ihr einen Stab gleichmäßig um $\Delta T$ erwärmen wollt, tragt
$T_{oben} = T_{unten} = \Delta T$ ein. Schreibt ihr nur $T_{oben} = \Delta T$ und lasst $h = 0$,
wirkt nur $\Delta T/2$ -- darauf weist euch eine gelbe Box hin.

Temperatur darf auch symbolisch sein (\code{dt}, \code{t1}, \code{a}). Damit könnt ihr genau die
Klausurfrage beantworten, bei der ein $\Delta T$ gesucht ist: erst symbolisch rechnen, dann im
CAS nach \code{dt} auflösen (Abschnitt~\ref{sec:export}).

\section{Symbolischer Modus}\label{sec:symbolisch}

Sobald in irgendeinem Feld ein Symbol steht (\code{q}, \code{F}, \code{EI}, \code{dt} \dots),
schaltet das Programm um. Oben rechts steht dann \emph{symb.} Das ist der Modus, in dem ihr für
\textbf{TM2} arbeiten wollt, weil dort fast nie Zahlen gegeben sind.

\begin{itemize}
  \item \textbf{Referenzlänge:} Steht in der Geometrie ein Symbol, ist das eure Referenzlänge.
        Steht dort keines, nimmt das Programm automatisch einen Rasterabstand als $l$ und gibt
        alle Stablängen als Vielfache davon aus (Raster 0,5~m, Stab 2~m $\rightarrow$ \code{4*l}).
        Ihr könnt ein System also ganz normal zeichnen und bekommt trotzdem eine symbolische
        Lösung.
  \item \textbf{Diagramme} werden symbolisch beschriftet. Die Zahlenfaktoren folgen dem
        Zahlenformat aus dem Obermenü (Seite~3): exakter Bruch, sonst -- bei \emph{Wurzel/Bruch}
        -- die Wurzel eines Bruchs, sonst eine Dezimalzahl. Eine 45°-Diagonale im Fachwerk steht
        damit als \code{50*sqrt(2)} da und nicht als Zufallsbruch. Schaltet ihr das Format um,
        wird neu gerechnet.
  \item \textbf{Exporte} (Abschnitt~\ref{sec:export}) sind vollständig symbolisch.
\end{itemize}

\textbf{Was nicht geht} (rote Box, die Berechnung bricht ab):

\begin{longtable}{p{0.45\linewidth}p{0.51\linewidth}}
  \toprule
  \textbf{Meldung} & \textbf{Grund} \\
  \midrule \endhead
  Nur ein Längensymbol erlaubt & Mehrere verschiedene Symbole in den Koordinaten. \\
  Koordinate kein Vielfaches von $l$ & Koordinaten wie \code{a+2}. \\
  Symbol \code{t} ist reserviert & \code{t} ist die Integrationsvariable des Exports, und der
    Rechner unterscheidet \code{t} und \code{T} nicht. Nehmt \code{dt}, \code{t1}, \dots \\
  Temperatur als Funktion von $x$ & $x$ in $T_{oben}$, $T_{unten}$, $\alpha_T$ oder $h$. \\
  Lastfunktion $f(x)$ & $x$ in $q$, $n$, $m$, $p_X$ oder $p_Y$ (numerisch ist das erlaubt). \\
  Bögen nicht symbolisch möglich & Bogenradius $\ne 0$. \\
  \bottomrule
\end{longtable}

Eine Mischung aus Symbol und Zahl bei den Steifigkeiten (z.\,B.\ Feder \code{3*EA*l} neben Stäben
mit $EA$ als Zahl) bricht \emph{nicht} ab: es wird gerechnet und exportiert. Ein gelber Hinweis
kommt nur dann, wenn der Zahlenwert das Ergebnis wirklich beeinflusst, also wenn der Stab
tatsächlich gebogen wird bzw.\ Normalkraft trägt (Abschnitt~\ref{sec:hinweise}). Sauber ist es
trotzdem, alle Steifigkeiten symbolisch zu machen.

Griechische Buchstaben erkennt der Rechner nicht als Symbol; schreibt \code{a} statt $\alpha$.

\bildplatz{tragwerk-symbolisch}{Symbolische Beschriftung der Schnittgrößen}

\section{Export ins CAS}\label{sec:export}

Hier wird es für \textbf{TM2} richtig praktisch: Das Programm rechnet nicht nur, sondern legt
euch die Ergebnisse als Variablen ins Dokument, mit denen ihr auf einer Calculator-Seite
weiterrechnen könnt. Die Exporte stehen im Obermenü (\taste{o}) auf Seite~2, und nach jedem
Export meldet sich eine gelbe Box mit den Namen der Variablen.

\begin{longtable}{p{0.32\linewidth}p{0.64\linewidth}}
  \toprule
  \textbf{Menüpunkt} & \textbf{Variablen im CAS} \\
  \midrule \endhead
  Alles exportieren & Systemmatrizen und Stabdaten (\code{sys\_k}, \code{sys\_f}, \code{sys\_u},
    \code{stab\_i}, \code{knot\_i}) \\
  Nur Biegelinien & \code{v\_EI\_stab\_i}: $EI\,w(x)$ je Stab aus der FEM-Lösung \\
  Nur Längslinien & \code{u\_EA\_stab\_i}: $EA\,u(x)$ je Stab \\
  Verschiebungslinien & beides zusammen \\
  KGV-Daten & \code{kgv\_delta}, \code{kgv\_delta0} (Abschnitt~\ref{sec:kgv}) \\
  Arbeitssatz & \code{w\_eq1}, \code{w\_eq2} (nur im PvV-Modus) \\
  Nur Schnittgrößen & \code{n\_stab\_i(x)}, \code{q\_stab\_i(x)}, \code{m\_stab\_i(x)} \\
  LGS-Matrix & \code{ggw\_a}, \code{ggw\_b}, \code{ggw\_x}, \code{lgs} \\
  Neue Biegelinie / Neue Längslinie & Rand-LGS, siehe unten \\
  \bottomrule
\end{longtable}

\subsection{LGS-Matrix: das Gleichgewicht zum Nachrechnen}

\emph{LGS-Matrix} schreibt das Gleichgewichtssystem heraus, das ihr in \textbf{Statik} und
\textbf{TM1} von Hand aufstellt: je Teilsystem $\sum H = 0$, $\sum V = 0$ und $\sum M = 0$ um
einen gut gewählten Punkt. In \code{lgs} stehen die Gleichungen als lesbarer Text, in
\code{ggw\_a}, \code{ggw\_b} und \code{ggw\_x} dieselben Gleichungen als Matrix, Rechte-Seite und
Unbekanntenvektor -- \code{simult(ggw\_a,ggw\_b)} löst sie also direkt.

Die Unbekannten heißen so, wie die Größen auch in der Explosionsansicht beschriftet sind:

\begin{longtable}{p{0.22\linewidth}p{0.74\linewidth}}
  \toprule
  \textbf{Name} & \textbf{Bedeutung} \\
  \midrule \endhead
  \code{A\_h}, \code{A\_v}, \code{M\_A} & Lagerkräfte und Einspannmoment (Lager werden mit
    Buchstaben durchnummeriert) \\
  \code{Gx\_i}, \code{Gy\_i} & Gelenkkraft am Momentengelenk~$i$, in globalen Komponenten \\
  \code{Gn\_i}, \code{Gq\_i} & bei einem N- oder Q-Gelenk: die Stabkraft längs bzw.\ quer zum Stab \\
  \code{Gm\_i} & das Moment, das ein N- oder Q-Gelenk überträgt \\
  \bottomrule
\end{longtable}

Der Hintergrund: Ein Momentengelenk überträgt $N$ und $Q$, aber kein Moment -- deshalb zwei
Kraftkomponenten. Ein \emph{N-Gelenk} überträgt dagegen $Q$ \emph{und} $M$, ein \emph{Q-Gelenk}
$N$ und $M$. Das Programm stellt genau diese Größen als Unbekannte ein, und zwar in
Stabrichtung, sauber ins globale System gedreht. Wenn ihr das LGS löst, kommen damit exakt die
Zahlen heraus, die auf Seite~1 der Explosionsansicht stehen.

\subsection{\texorpdfstring{Rand-LGS: \code{wneu}, \code{uneu}, \code{wv}, \code{uv}, \code{randbed}}{Rand-LGS: wneu, uneu, wv, uv, randbed}}

Die Menüpunkte \emph{Neue Biegelinie} und \emph{Neue Längslinie} rufen dieselbe Rechnung auf, und
zwar die, die ihr in \textbf{TM2} von Hand macht: Für jeden Stab wird angesetzt

\[
  EI\,w_i(x) = w_{p,i}(x) + c_1\frac{x^3}{6} + c_2\frac{x^2}{2} + c_3 x + c_4,
  \qquad
  EA\,u_i(x) = u_{p,i}(x) + C_1 x + C_2
\]

und die Konstanten werden aus \emph{allen} Rand- und Übergangsbedingungen des Systems bestimmt.
Exportiert werden:

\begin{longtable}{p{0.22\linewidth}p{0.74\linewidth}}
  \toprule
  \textbf{Variable} & \textbf{Inhalt} \\
  \midrule \endhead
  \code{wneu\_i(x)} & $EI\,w_i(x)$, also die Biegelinie in der Form der Handrechnung \\
  \code{uneu\_i(x)} & $EA\,u_i(x)$; die Ableitung liefert die Normalkraft \\
  \code{wv\_i(x)} & $w_i(x)$, die Durchbiegung selbst \\
  \code{uv\_i(x)} & $u_i(x)$, die Längsverschiebung selbst \\
  \code{randbed} & Spaltenvektor mit allen verwendeten Bedingungen als Text \\
  \code{randinfo} & Legende: starre Stäbe, verwendete Methode, Bedeutung der Funktionen \\
  \bottomrule
\end{longtable}

In \code{randbed} stehen die Bedingungen so, wie man sie von Hand aufschreibt, mit den
Schnittgrößen als Ableitungen, also \code{-EIw'''} statt $Q$, \code{-EIw''} statt $M$ und
\code{EAu'} statt $N$:

\begin{center}
\code{["w1(x1=0)=0";"EIw1''(x1=0)=0";"w1(x1=l)=w2(x2=0)";\dots]}
\end{center}

Was vor einer Bedingung schon feststeht, steht dort gleich als Wert -- genau wie bei euch auf dem
Papier. Ist Stab~2 dehnstarr und gehalten, wird aus \code{"w1(x1=l)=u2(x2=3*l)"} also
\code{"w1(x1=l)=0"}. Und wenn sich eine der beiden Linien ohne die andere bestimmen lässt --
typisch, wenn dehnstarre oder biegesteife Stäbe die Kopplung auflösen --, ist \code{randbed} mit
Überschriften in zwei Gruppen geteilt:

\begin{center}
\code{["--- nur Biegelinie w ---";\dots;"--- nur Laengslinie u ---";\dots]}
\end{center}

Für die Biegelinie genügen dann die Bedingungen der ersten Gruppe. Das ist genau der Fall, den
Klausuraufgaben gerne bauen: Es wird nur die Biegelinie gefragt, und die Längslinie braucht man
gar nicht erst. Bei gekoppelten Systemen (Rahmen, deren Normalkräfte die Querkräfte des
Nachbarstabs aufnehmen) bleibt die Liste ungeteilt, weil sich das dort nicht trennen lässt.

\subsection{Nach einer gesuchten Größe auflösen}

Weil \code{wv} und \code{uv} die Verschiebungen selbst sind, könnt ihr die typischen Fragen
direkt auf dem Rechner beantworten:

\begin{itemize}
  \item \enquote{Bestimmen Sie $\Delta T$, sodass sich Punkt A nicht verschiebt:}\\
        \code{solve(uv1(l)=0,dt)}
  \item \enquote{Welche Kraft $F$ verschiebt den Punkt um 0,5?}\\
        \code{solve(wv2(2*l)=0.5,f)}
\end{itemize}

Voraussetzung ist nur, dass die gesuchte Größe vorher als Symbol eingegeben wurde. Die Funktionen
sind exakt; wenn ihr Dezimalzahlen wollt, nehmt \code{approx(wneu1(x))}.

\textbf{Grenzen des Rand-LGS:} keine Bögen, kein KGV-Einheitszustand (Zustand~0 wählen), keine
Lasten in frei beweglichen Richtungen. Wenn etwas schiefgeht, steht der Grund in einer roten Box
und zusätzlich in \code{randbed} als \code{["Fehler: \dots"]}; die alten Funktionen werden vorher
gelöscht, es bleibt also nichts Veraltetes stehen, mit dem ihr aus Versehen weiterrechnet.

\section{Kraftgrößenverfahren (KGV)}\label{sec:kgv}

Mit \emph{Auto-KGV} (Obermenü Seite~1) sucht sich das Programm selbst ein statisch bestimmtes
Hauptsystem und zeigt oben rechts den Grad der statischen Unbestimmtheit $n$. Mit den
Zifferntasten \taste{1}\,\dots\,\taste{9} schaltet ihr zwischen dem Lastzustand (Zustand~0) und
den Einheitszuständen $X_i = 1$ um -- und seht dabei jedes Mal die kompletten Schnittgrößen.

Gelöst wird in dieser Reihenfolge: Einspannmomente, Lagerkräfte, Momentengelenke und zuletzt
unbelastete Pendelstäbe. Der Export \emph{KGV-Daten} liefert $\delta_{ij}$ (\code{kgv\_delta}) und
$\delta_{i0}$ (\code{kgv\_delta0}); daraus folgt auf dem Rechner
$\underline{X} = -\delta^{-1}\,\delta_0$.

\textbf{Wichtig bei gelösten Pendelstäben:} Der Stab ist im Hauptsystem entfernt, die
Einheitskräfte greifen an seinen Knoten an. Seine eigene Längenänderung $NL/(EA)$ steckt deshalb
nicht in den Knotenverschiebungen und wird zu $\delta_{ii}$ addiert -- aber nur zu diesem einen
Eintrag. Der Grund: Im Einheitszustand $i$ trägt allein der gelöste Stab~$i$ die Kraft $N = 1$,
alle anderen gelösten Stäbe sind ebenfalls aufgetrennt und damit kraftfrei. Dehnsteife Stäbe
liefern keinen Beitrag, und bei gelösten Lagern oder eingefügten Gelenken entfällt der Zusatz,
weil der Stab dort im System bleibt.

\section{Prinzip der virtuellen Verschiebungen (PvV)}

\taste{v} öffnet den Verschiebungszustand. Solange das System nicht kinematisch ist, kommt die
Meldung \emph{System ist starr! [Enter] für PvV-Modus}. Nach \taste{Enter} wählt ihr mit
\taste{Enter} auf einem Knoten oder Stab, welche Bindung ihr lösen wollt:

\begin{itemize}
  \item am Knoten: Lager X, Lager Y, Lager M entfernen oder ein Gelenk einfügen (letzteres nur
        bei genau zwei Stäben),
  \item am Stab: M-, N- oder Q-Gelenk an Anfang oder Ende.
\end{itemize}

Danach zeigt euch das Programm die kinematische Kette, auf Wunsch den Polplan (\taste{p}) und in
den Menüs die Starrkörperverschiebungen, die Lastangriffspunkte und die einzelnen Arbeitsanteile
$W$. Der Menüpunkt \emph{Arbeitssatz} exportiert die fertige Gleichung als \code{w\_eq1} und
\code{w\_eq2}. Für \textbf{Statik} ist das die schnellste Kontrolle, ob eure Handrechnung stimmt.

\bildplatz{tragwerk-pvv}{PvV-Modus mit Polplan und Arbeitsanteilen}

\section{Meldungen im Tragwerksskript}

Alle Meldungen erscheinen als Box in der Bildmitte. \taste{Esc} schließt immer die oberste Box:
zuerst Eingabe oder Menü, dann die Box, dann Sondermodi und die Auswahl.

\subsection{Warnungen (rot bzw.\ orange)}

\begin{longtable}{p{0.34\linewidth}p{0.29\linewidth}p{0.31\linewidth}}
  \toprule
  \textbf{Meldung} & \textbf{Wann} & \textbf{Warum} \\
  \midrule \endhead
  Symbolischer Modus nicht möglich & Symbolische Eingabe, die das Verfahren nicht abbilden kann
    (Liste in Abschnitt~\ref{sec:symbolisch}) & Verhindert falsche Ergebnisse; Berechnung und
    Rand-LGS bleiben gesperrt, auch nachdem ihr die Box geschlossen habt \\
  System ist kinematisch! & Das Gleichungssystem ist singulär & Das System ist ein Mechanismus,
    es gibt keine eindeutige Lösung \\
  Starrmodus \& statisch unbestimmt! & Statisch unbestimmtes System mit sehr steifen Stäben
    ($EA$, $EI \ge 10^7$, also auch mit den Standardwerten) & In unbestimmten Systemen hängt die
    Kraftverteilung von den Steifigkeitsverhältnissen ab -- mit Standardwerten rechnet ihr dann
    etwas aus, das mit der Aufgabe nichts zu tun hat \\
  System ist starr! [Enter] für PvV-Modus & PvV-Ansicht bei nicht kinematischem System &
    Für PvV muss erst eine Bindung gelöst werden \\
  Wirklich alles loeschen? & \taste{c} bzw.\ \taste{Clear} & Sicherheitsfrage: \taste{Enter}
    löscht, \taste{Esc} bricht ab \\
  \bottomrule
\end{longtable}

\subsection{Hinweise (gelb)}\label{sec:hinweise}

Diese Box kommt nach der Berechnung; das Ergebnis ist trotzdem gültig. Nach \taste{Esc} kommt sie
erst wieder, wenn sich am System etwas geändert hat -- sie nervt euch also nicht bei jedem
Rechnen erneut.

\begin{longtable}{p{0.34\linewidth}p{0.29\linewidth}p{0.31\linewidth}}
  \toprule
  \textbf{Hinweis} & \textbf{Wann} & \textbf{Warum} \\
  \midrule \endhead
  $T_{oben} \ne T_{unten}$, aber $h = 0$ & Temperaturunterschied ohne Höhe &
    Es wirkt nur die mittlere Erwärmung, keine Krümmung \\
  Temperatur symbolisch: Diagrammwerte ohne Temperaturanteil & Symbolischer Modus mit Temperatur &
    Die Diagrammbeschriftung entsteht durch Superposition und enthält die Temperatur nicht; der
    Export schon \\
  Stab~$i$ wird auf Biegung belastet, $EI$ ist dort eine Zahl & Symbolischer Modus, der Stab wird
    wirklich gebogen (kein Pendelstab) und hat ein Zahlen-$EI$, während andere Steifigkeiten
    symbolisch sind & Nur dann beeinflusst der Zahlenwert das Ergebnis. Statisch bestimmt: die
    Schnittgrößen bleiben richtig, Biegelinie und Export mischen Zahl und Symbol. Statisch
    unbestimmt: auch die Schnittgrößen können falsch werden (Symbole rechnen intern mit 1).
    Abhilfe: $EI$ symbolisch, biegesteif oder Gelenke an beiden Enden \\
  Stab~$i$ trägt Normalkraft, $EA$ ist dort eine Zahl & wie oben mit $N \ne 0$ & Abhilfe:
    dehnsteif, Standard-EA dehnstarr oder $EA$ symbolisch \\
  \bottomrule
\end{longtable}

\subsection{Export- und Rand-LGS-Meldungen}

Erfolgreiche Exporte melden sich mit dem Titel \emph{Export}, Fehlschläge mit \emph{Export
fehlgeschlagen}. Typische Fehlermeldungen: \emph{Bitte zuerst das System berechnen}, \emph{KGV
nicht aktiv oder System statisch bestimmt}, \emph{bitte zuerst in den PvV-Modus wechseln},
\emph{Dieser Export erfordert das CAS}. Beim Rand-LGS zusätzlich: nicht quadratisch, singulär,
Bögen nicht unterstützt, Last in frei beweglicher Richtung, unendliche Zwangskraft bei
behinderter Temperaturverformung eines starren Stabes.

\section{Was das Tragwerksskript nicht kann}

\begin{itemize}
  \item Nur ebene Stabtragwerke -- keine räumlichen Systeme, keine Flächentragwerke.
  \item Lineare Theorie I.~Ordnung: keine Stabilität (Knicken), keine großen Verformungen, keine
        Plastizität, kein Schubverformungsanteil.
  \item Keine Wanderlasten, keine Einflusslinien, keine Lastfallkombinationen.
  \item Der symbolische Modus deckt nur die in Abschnitt~\ref{sec:symbolisch} genannten Fälle ab.
  \item Bögen werden für die Diagramme in Teilstäbe zerlegt; im Rand-LGS-Export gehen sie nicht.
  \item Das automatisch gewählte KGV-Hauptsystem muss nicht das sein, das ihr von Hand wählen
        würdet -- die Ergebnisse stimmen trotzdem, die Zwischenwerte $\delta_{ij}$ aber natürlich
        nur zu \emph{diesem} Hauptsystem.
  \item Es gibt keine Materialdatenbank und keine Bemessung; $EA$ und $EI$ gebt ihr ein.
\end{itemize}

\section{Wie gerechnet wird}

Das braucht ihr nicht zu wissen, um das Programm zu benutzen -- aber es erklärt, warum manche
Eingabe so ist, wie sie ist.

\textbf{Diagramme und Auflagerkräfte} entstehen mit dem Weggrößenverfahren (also finite Elemente
für den ebenen Stab): Für jeden Stab wird die Steifigkeitsmatrix aufgestellt, Gelenke werden
kondensiert, Lasten als Festeinspannkräfte eingerechnet, Federn und gedrehte Lager transformiert.
Aus den Knotenverschiebungen folgen die Stabendkräfte und daraus die Verläufe. Deshalb braucht
jeder Bereich, den ihr getrennt betrachten würdet, einen eigenen Stab. Starre Stäbe werden nicht
mit großen Zahlen angenähert, sondern als exakte Zwangsbedingungen geführt; ihre Schnittgrößen
kommen aus den zugehörigen Zwangskräften.

\textbf{Die exportierten Linien} entstehen unabhängig davon über das Rand-LGS: Ansatz je Stab,
alle Rand- und Übergangsbedingungen, Lösung im CAS mit \code{simult}. Vorher wertet das Programm
alle Bedingungen aus, die nur eine Unbekannte enthalten, streicht Unbekannte, die nur in einer
einzigen Zeile vorkommen, und setzt reine Zahlenbeziehungen ein; was übrig bleibt, wird in
unabhängige Blöcke zerlegt und blockweise gelöst. Dadurch sind die Funktionen exakt und im
symbolischen Modus auch symbolisch. Für starre Stäbe gibt es zwei Wege: direkt
(Starrkörperbewegung plus Gleichgewicht) oder als Grenzwert $EA, EI \to \infty$.

\textbf{Die symbolische Beschriftung der Diagramme} entsteht durch Superposition: Jede Last wird
einzeln mit einem Platzhalterwert gerechnet und das Ergebnis als Vielfaches des Symbols und der
Referenzlänge wieder zusammengesetzt. Temperaturlasten passen nicht in dieses Schema -- daher der
Hinweis oben.

%=====================================================================
\part{Querschnittsskript}

\section{Ein erster Querschnitt}

Das Querschnittsmodul ist vor allem für \textbf{TM2} gedacht: Querschnittswerte, Spannungen,
Schub, Torsion, Kernfläche. Für \textbf{TM1} könnt ihr es für Schwerpunkt- und
Flächenträgheitsmomente benutzen.

\begin{enumerate}
  \item \taste{m} öffnet das Menü für massive Formen, \taste{d} das für dünnwandige Profile.
  \item Form auswählen (Ziffer) und die nötigen Punkte klicken.
  \item Mit dem Zeiger auf einem Element und \taste{Enter} öffnet sich dessen Bearbeitungsmenü
        (Koordinaten, Dicke, Loch/Positiv, Löschen); ein Klick vorher ist nicht nötig.
  \item \taste{e} zeigt die Querschnittswerte, \taste{t} die Einzelwerte je Element.
  \item \taste{b} öffnet das Berechnungsmenü für Spannungen, Schub, Kernfläche und Verwölbung.
\end{enumerate}

Zeichnen könnt ihr grob -- die Koordinaten korrigiert ihr danach im Elementmenü exakt. Das ist
meist schneller, als von Anfang an genau zu klicken.

\bildplatz{querschnitt-eingabe}{Eingabe eines dünnwandigen Profils}

\section{Elemente zeichnen}

\begin{longtable}{p{0.30\linewidth}p{0.66\linewidth}}
  \toprule
  \textbf{Massiv (\taste{m})} & \textbf{Klicks} \\
  \midrule \endhead
  Rechteck & 2 (gegenüberliegende Ecken) \\
  Kreis & 2 (Mittelpunkt, Randpunkt) \\
  Dreieck & 3 \\
  Trapez/Polygon & beliebig viele, \taste{Enter} schließt \\
  Kreisausschnitt & 3 (Mittelpunkt, Start, Ende) \\
  Kreisabschnitt & 3 \\
  Dicke Linie & 2 (Rechteck um die Verbindungslinie, Dicke $t$) \\
  Kraft platzieren & 1 \\
  \bottomrule
\end{longtable}

\begin{longtable}{p{0.30\linewidth}p{0.66\linewidth}}
  \toprule
  \textbf{Dünnwandig (\taste{d})} & \textbf{Klicks} \\
  \midrule \endhead
  Linienzug & beliebig viele, \taste{Enter} schließt ab \\
  Kreis & 2 \\
  Kreisbogen & 3 \\
  Kraft platzieren & 1 \\
  \bottomrule
\end{longtable}

Geklickte Punkte rasten auf dem Raster ein -- und auf den \textbf{Punkten der vorhandenen
Elemente}: Ecken von Rechtecken (auch die beiden, die ihr nicht geklickt habt), Dreiecken und
Polygonen, Mittel- und Randpunkte von Kreisen und Bögen, Endpunkte der Linien, dazu die schon
geklickten Punkte des Elements, das gerade entsteht. Genommen wird immer der nächste Punkt. Damit
trefft ihr auch Punkte, die nach einer exakten Koordinateneingabe oder einer KOS-Verschiebung
nicht mehr auf dem Raster liegen.

Die Wanddicke kommt aus der Standarddicke $t$ (Optionsmenü \taste{o}); einzelne Elemente ändert
ihr nachträglich über Hover~+~\taste{Enter}. \taste{Backspace} löscht das letzte Element,
\taste{c} löscht alles -- nach einer Rückfrage, die ihr mit \taste{Enter} bestätigt oder mit
\taste{Esc} abbrecht.

\textbf{Überlappung:} Überdeckt ein neues Massivelement ein vorhandenes wirklich (gemeinsame
Kanten zählen nicht), fragt das Programm nach: \taste{1} zieht es als Loch ab, \taste{2} addiert
es numerisch (Rasternäherung, im Ergebnisfeld steht dann \emph{Warnung: Numerisch berechnet!}),
\taste{3} verwirft die Eingabe. Für Löcher ist \taste{1} der saubere Weg, weil dann weiter mit
exakten Formeln gerechnet wird.

\section{Koordinatensystem, Einheiten, Optionen}

Im Optionsmenü (\taste{o}):

\begin{itemize}
  \item Raster und Standarddicke $t$,
  \item KOS-Ebene und KOS drehen (in Schritten von 90°), KOS-Größe,
  \item Zahlenformat (Wurzel/Bruch, Bruch, Dezimal),
  \item Einheiten für Kraft, Länge und Moment,
  \item Profilbeiwert $\xi$ für die Torsion offener Profile,
  \item auf Seite~2 (\emph{Weitere Optionen}, Pfeil links/rechts): KOS in den Schwerpunkt
        verschieben, KOS entlang $y$ bzw.\ $z$ um einen Betrag verschieben und \emph{FTM-Bezug}.
\end{itemize}

\textbf{KOS entlang einer Achse verschieben:} Ihr gebt den Betrag in der aktuellen Längeneinheit
ein, positiv heißt in positive Achsrichtung. Der Querschnitt bleibt auf dem Bildschirm stehen,
das KOS (und mit ihm das Raster) wandert. Zusammen mit dem FTM-Bezug \emph{KOS} bekommt ihr so
die Trägheitsmomente um jede beliebige parallele Achse.

\textbf{FTM-Bezug} legt fest, worauf sich die angezeigten Flächenträgheitsmomente beziehen: auf
den \emph{Schwerpunkt} (Standard, Anzeige $I_{y,S}$) oder auf den \emph{Ursprung des gezeichneten
KOS} (Anzeige $I_y$~(KOS)). Beides braucht man: Mit dem Bezug \emph{KOS} seht ihr in der
Einzelwerte-Tabelle die Steiner-Anteile zum Ursprung, so wie man eine Handrechnung aufbaut; mit
dem Bezug \emph{Schwerpunkt} bekommt ihr direkt die Werte, mit denen ihr Spannungen rechnet.

Alle Ein- und Ausgaben beziehen sich auf das \emph{angezeigte} Koordinatensystem: Dreht ihr das
KOS, drehen sich Koordinaten, Trägheitsmomente, eingegebene Schnittgrößen und Ergebnisse mit.
Das ist praktisch, wenn die Aufgabe das KOS anders herum definiert als ihr gezeichnet habt.

\section{Querschnittswerte}

\taste{e} zeigt das Ergebnisfeld: Fläche, Schwerpunkt, Flächenmomente, Hauptträgheitsmomente mit
Hauptachsenwinkel, Widerstandsmomente und -- nach einer Schubrechnung -- die Schubwerte.
\taste{t} zeigt dieselben Werte je Element, inklusive der Steiner-Anteile; das ist die Tabelle,
die ihr in der Klausur sowieso aufschreibt.

Zwei Zeilen im Ergebnisfeld solltet ihr immer anschauen:

\begin{itemize}
  \item \textcolor{green!50!black}{\textbf{Hauptachsen parallel zu den Achsen ($I_{yz} = 0$)}}:
        Die einfachen Formeln gelten.
  \item \textcolor{orange!80!black}{\textbf{Hauptachsen nicht parallel zu den Achsen}}:
        schiefe Biegung; $\sigma_x$ und der Schubfluss brauchen die gekoppelten Formeln. Das
        Programm rechnet damit, ihr müsst es aber wissen, wenn ihr die Rechnung von Hand
        nachvollziehen wollt.
\end{itemize}

\bildplatz{querschnitt-werte}{Ergebnisfeld mit Querschnittswerten}

\section{\texorpdfstring{Normalspannung $\sigma_x$}{Normalspannung sigma-x}}

Im Berechnungsmenü (\taste{b}) unter \emph{$\sigma_x$-Berechnung}:

\begin{itemize}
  \item \emph{Normalspannung} $(N, M_y)$ für gerade Biegung,
  \item \emph{Schiefe Spannung} $(N, M_y, M_z)$,
  \item \emph{$\sigma_x$ aus Kräften}, wenn ihr Kräfte platziert habt (mit Querkräften wird der
        Abstand $\Delta x$ gefragt, siehe unten).
\end{itemize}

Ihr bekommt Minimum und Maximum mit ihren Orten, die Nulllinie und die Spannungsverteilung über
den Querschnitt. Gerechnet wird immer mit der allgemeinen Biegeformel inklusive
$I_{yz}$-Kopplung, also auch für unsymmetrische Querschnitte richtig.

\textbf{Ergebnistabelle:} Mit \taste{t} kommt ihr zur Tabelle, die Pfeiltasten hoch/runter
scrollen durch sie. Außer den Schnittgrößen und
Extremwerten steht dort:

\begin{itemize}
  \item das \textbf{Hauptachsensystem}: Winkel $\varphi^*$ von $y$ nach $\eta$, $I_\eta = I_1$,
        $I_\zeta = I_2$ und die \textbf{Momente um die Hauptachsen}
        $M_\eta = M_y\cos\varphi + M_z\sin\varphi$, $M_\zeta = -M_y\sin\varphi + M_z\cos\varphi$,
  \item die \textbf{Biegenormalspannung} in MPa als Gleichung, einmal im gezeigten KOS
        ($\sigma_x = c_0 + c_y\,y + c_z\,z$, $y$ und $z$ vom KOS-Ursprung aus) und einmal im HAS
        ($\sigma_x = N/A + (M_\eta/I_\eta)\,\zeta - (M_\zeta/I_\zeta)\,\eta$, als Zahlen),
  \item die \textbf{neutrale Faser} $\sigma_x = 0$ in beiden Systemen, z.\,B.\ $z = m\,y + n$.
\end{itemize}

Liegt euer KOS nicht im Schwerpunkt, steckt das in der Konstanten $c_0$ der KOS-Gleichung; die
Lage des Schwerpunkts steht darunter. Die HAS-Gleichung ist immer auf den Schwerpunkt bezogen.

\textbf{Querkräfte mit Hebelarm:} Habt ihr Kräfte eingezeichnet, die in der Querschnittsebene
wirken, fragt die Maske zusätzlich $\Delta x$ ab -- den Abstand zwischen dem Schnitt, den ihr
anschaut, und der Ebene, in der die Kräfte angreifen (im $yz$-KOS also entlang $x$). Die
Querkräfte liefern dann

\[
  M_y = -\Delta x\,F_z, \qquad M_z = \Delta x\,F_y ,
\]

und das wird mit euren eingegebenen Momenten überlagert. Wichtig für die Vorzeichen: Es zählt
nur der \emph{Abstand}. Ob die Kräfte vor oder hinter dem Schnitt sitzen, ist für diesen Anteil
egal -- rechnet man es für beide Seiten mit $\vec M = \pm\,\vec r \times \vec F$ nach, kommt
jedes Mal $M_y = -a F_z$ heraus. Denkt an einen Kragarm: Eine Last nach unten macht immer Zug
oben, egal ob der Kragarm nach links oder nach rechts auskragt. Ein negatives $\Delta x$ nimmt
das Programm deshalb als Betrag. Bei \emph{$\sigma_x$ aus Kräften} wird nur $\Delta x$ gefragt.

\section{Schubspannung und Schubmittelpunkt}

\emph{Schubspannung} im Berechnungsmenü fragt $Q_y$ und $Q_z$ ab (im angezeigten KOS); platzierte
Kräfte werden dazugezählt. Danach wählt ihr mit \taste{v} den Verlauf, die Ziffern schalten um:

\begin{longtable}{p{0.10\linewidth}p{0.86\linewidth}}
  \toprule
  \textbf{Taste} & \textbf{Anzeige} \\
  \midrule \endhead
  1, 2 & Momentendichten $y\,h(y)$ bzw.\ $z\,h(z)$ \\
  3, 4 & statische Momente $S_y$, $S_z$ \\
  5, 6 & Schubspannungsanteile $\tau_y$, $\tau_z$ \\
  7 & Resultierende aus Querkraft \\
  8 & Torsionsanteil \\
  9 & Gesamtschubspannung \\
  0 & Laufvariable $s$ \\
  \bottomrule
\end{longtable}

Der \emph{Schubmittelpunkt} wird aus den Momenten der Schubflüsse für Einheitsquerkräfte
bestimmt; mit der Momententabelle (\taste{t} im Zusatzmenü) könnt ihr die Rechnung Segment für
Segment nachvollziehen.

\textbf{Berechnet wird:}

\begin{itemize}
  \item offene dünnwandige Profile, auch unsymmetrische (mit $I_{yz}$-Kopplung),
  \item Profile mit genau einer geschlossenen Zelle, auch mit offenen Ästen -- mit Symmetrieachse
        parallel zur Querkraft über die Symmetrie, sonst über die Verträglichkeit (siehe unten),
  \item massive Querschnitte mit $\tau = QS/(Ib)$, aber nur bei $I_{yz} = 0$.
\end{itemize}

Nicht berechnet werden: mehrzellige Profile, gemischt massiv-dünnwandige Querschnitte und massive
Querschnitte mit gedrehten Hauptachsen. Dann kommt eine rote Box mit dem Grund.

\subsection{Geschlossene Zellen: warum die Symmetrie so wichtig ist}

Bei einer geschlossenen Zelle bekommt ihr den Schubfluss nicht allein aus dem Gleichgewicht:
Schneidet man die Zelle auf, fehlt ein konstanter Umlaufanteil $q_0$. In \textbf{TM2} bestimmt
man den aus der Symmetrie -- bei einer Querkraft parallel zur Symmetrieachse ist der Schubfluss
antisymmetrisch und damit genau \emph{auf} der Achse null.

Genau das macht das Programm:

\begin{itemize}
  \item Es prüft die Spiegelsymmetrie, indem es die Mittellinien abtastet. Die Symmetrie wird
        also auch erkannt, wenn ihr die beiden Hälften unterschiedlich in Segmente geteilt habt.
  \item Das Ergebnisfeld zeigt bei geschlossenen Profilen \code{Zellsymmetrie: <Achse>} bzw.\
        \code{keine erkannt (q0 offen)} -- und zwar schon vor der Schubrechnung.
  \item Die Zelle wird \emph{auf der Symmetrieachse} aufgeschnitten. Die Laufvariable beginnt
        damit dort, wo $S = 0$ ist; die Marke im Verlauf heißt \emph{Start (Symmetrieachse)}.
  \item Auch die angezeigten $S$-Verläufe sind auf diesen Schnitt bezogen, nicht nur $\tau$.
\end{itemize}

\subsection{Ohne Symmetrieachse: die Verträglichkeit}

Fehlt die Symmetrieachse parallel zur Querkraft -- unsymmetrische Zelle, Querkraft quer zur
einzigen Achse, oder das Programm erkennt die Symmetrie aus irgendeinem Grund nicht --, dann hilft
die Symmetrie nicht weiter. Das Programm schneidet die Zelle trotzdem auf, rechnet den offenen
Verlauf und holt sich $q_0$ aus der \textbf{Verträglichkeit}: Unter einer Querkraft im
Schubmittelpunkt verdrillt sich die Zelle nicht, also

\[
  \oint \frac{q}{G\,t}\,ds = 0 .
\]

Das ist genau die Handrechnung, die ihr sonst für die statisch unbestimmte Zelle machen würdet.
Praktisch daran: Die Bedingung zerfällt in $\oint S_z/t\,ds = 0$ und $\oint S_y/t\,ds = 0$, das
Programm verschiebt also jedes statische Moment um seine eigene Konstante. Damit stimmen nicht nur
$\tau$, sondern auch die angezeigten $S$-Verläufe, und der Schubmittelpunkt ist für jede einzellige
Zelle bestimmt -- also auch \emph{Torsion aus Kräften}. Beim symmetrischen Profil kommt exakt die
antisymmetrische Lösung heraus; die Marke im Verlauf heißt dann \emph{Schnitt (q0 aus
Verträglichkeit)}.

\textbf{Die Warnung bleibt trotzdem}, und zwar an zwei Stellen: bei der Schubrechnung selbst
(die Symmetrie fehlt, $q_0$ kam aus der Verträglichkeit) und wenn ihr einen Verlauf anschaut,
dessen Richtung keine Symmetrieachse hat. Habt ihr z.\,B.\ nur die $z$-Achse als Symmetrieachse,
dann kommen $S_y$ und $\tau_z$ aus der Symmetrie, $S_z$ und $\tau_y$ aus der Verträglichkeit --
genau dann kommt die Warnung, und sie sagt dazu, dass $S$ und $\tau$ trotzdem richtig sind. Ihr
sollt wissen, dass ihr in der Klausur an dieser Stelle die Verträglichkeitsbedingung braucht.

Bei $I_{yz} \ne 0$ und dünnwandigem Profil erscheint zusätzlich ein gelber Hinweis mit dem
Hauptachsenwinkel: Der Schubfluss wurde mit der Kopplung

\[
  \tau\,t = -\frac{Q_z\,(I_z S_y + I_{yz} S_z) + Q_y\,(I_y S_z + I_{yz} S_y)}{I_y I_z - I_{yz}^2}
\]

berechnet und nicht mit $\tau = QS/(It)$.

\bildplatz{querschnitt-schub}{Schubspannungsverlauf mit Schubmittelpunkt}

\section{Torsion}

Torsion rechnet das Programm nur für rein dünnwandige Profile:

\begin{itemize}
  \item offen: $I_T = \sum \frac{\xi}{3} h t^3$, \quad $\tau_{max} = \dfrac{M_T\,t_{max}}{I_T}$,
  \item eine geschlossene Zelle (Bredt): $I_T = \dfrac{4A_m^2}{\oint \frac{ds}{t}}$,
        \quad $\tau_{max} = \dfrac{M_T}{2 A_m t_{min}}$.
\end{itemize}

Das Torsionsmoment könnt ihr eingeben oder aus platzierten Kräften bestimmen lassen; Bezugspunkt
ist dann der \emph{Schubmittelpunkt}. Für massive und gemischte Querschnitte gibt es keine
TM2-Formel, für mehrzellige Profile ebenfalls nicht: Das Programm lehnt mit Begründung ab und
rechnet bewusst keine Prandtl-Lösung, die ihr in der Klausur ohnehin nicht verwenden dürftet.

\section{Verwölbung}

Im Berechnungsmenü unter \emph{7. Verwölbung} bekommt ihr die Verwölbung dünnwandiger Profile
nach Kapitel~6 der Formelsammlung:

\[
  u_x(s) = \int_s \left[ \frac{M_T}{2\,G\,A_m\,h(s)} - r_\perp\,\vartheta \right] ds + c,
  \qquad \vartheta = \frac{M_T}{G\,I_T}, \qquad I_T = \frac{4A_m^2}{\oint ds/h}.
\]

Ein Dialog fragt $M_T$ (vorbelegt mit dem Moment eurer platzierten Kräfte um den
Schubmittelpunkt, sonst leer) und den Schubmodul $G$ in N/mm$^2$ (vorbelegt 81\,000). Das
$M_T$ gilt \textbf{nur in dieser Ansicht} -- \taste{Esc} schließt sie und verwirft es, an die
Torsion der Schubspannungsrechnung wird es nicht übergeben. Der Verlauf wird wie die
Schubverläufe normal zur Wand gezeichnet, $M_T$ als Drehpfeil um den Pol; \taste{e} zeigt
$\vartheta$, $I_T$, $A_m$ und $u_{x,max}$.

Drei Dinge, die ihr aus der Handrechnung kennt und die das Programm genauso macht:

\begin{itemize}
  \item \textbf{Pol ist der Schubmittelpunkt}, denn um den dreht sich der Querschnitt bei
        Torsion. Bei doppelt symmetrischen Profilen ist das der Schwerpunkt.
  \item $r_\perp$ \textbf{hat ein Vorzeichen}: positiv, wenn der Fahrstrahl vom Pol zum
        Laufpunkt beim Fortschreiten im Drehsinn des positiven Torsionsmoments dreht. Bei einer
        Zelle mit dem Pol innen ist das überall der positive Abstand -- deshalb funktioniert die
        Handrechnung mit \enquote{Abstand} so gut. Beim C-Profil liegt der Pol außerhalb, und der
        Stegbeitrag geht mit anderem Vorzeichen ein als die Flanschbeiträge.
  \item \textbf{Die Konstante $c$} legt das Nullniveau fest: $\int u_x\,h\,ds = 0$, also keine
        mittlere Längsverschiebung. Bei symmetrischen Profilen ist $u_x$ damit auf der Achse null
        -- genau das bekommt ihr, wenn ihr dort mit $c = 0$ startet.
\end{itemize}

\textbf{Offene Profile} rechnet das Programm auch: Der Bredt-Term entfällt, es bleibt
$u_x = -\vartheta \int r_\perp\,ds + c$. Das ist Standardtheorie, steht aber nicht in der
Formelsammlung -- deshalb kommt eine Warnung. Ein Quadratrohr wölbt sich nicht; das Rechteckrohr
$b \times h$ hat an den Ecken $|u_x| = \vartheta\,b\,h\,|b-h|/(4(b+h))$, abwechselnd nach vorn und
hinten. Mehrzellige, massive und gemischte Querschnitte: Fehler mit Begründung.

\section{Kernfläche}

\taste{k} schaltet den Kernflächenmodus weiter: umhüllende Geraden, Kernfläche, aus. Die
Kernfläche ist der Bereich, in dem eine Normalkraft angreifen darf, ohne dass im Querschnitt ein
Vorzeichenwechsel der Spannung auftritt -- also die klassische Frage \enquote{wo darf die Last
stehen, damit nichts auf Zug kommt}.

\section{Meldungen im Querschnittsskript}

Meldungen erscheinen als Box in der Bildmitte: rot, wenn etwas nicht berechnet wurde oder eine
Eingabe ungültig war, gelb bei Hinweisen. \taste{Esc} schließt zuerst die Box; sie verschwindet
auch, wenn ihr neu rechnet oder den Querschnitt ändert.

\begin{longtable}{p{0.40\linewidth}p{0.55\linewidth}}
  \toprule
  \textbf{Meldung} & \textbf{Bedeutung} \\
  \midrule \endhead
  Schub: Kombination massiv und dünnwandig & Für gemischte Querschnitte gibt es keine TM2-Formel \\
  Schub: Hauptachsen nicht parallel (massiv) & $\tau = QS/(Ib)$ gilt nur bei $I_{yz} = 0$ \\
  Schub: nur für gerade dünnwandige Elemente & Dünnwandige Kreise und Bögen sind nicht im
    Profilgraph \\
  Schub: mehrzelliges Profil & Statisch unbestimmt \\
  Warnung: Querkraft quer zur Symmetrieachse & $q_0$ kommt aus der Verträglichkeit, $S$ und
    $\tau$ sind trotzdem richtig; die Meldung nennt die erkannten Achsen \\
  Warnung: Geschlossene Zelle ohne erkannte Symmetrie & $q_0$ aus der Verträglichkeit, $S$ und
    $\tau$ trotzdem richtig \\
  Warnung: $S$ bzw.\ $\tau$ in der Richtung ohne Symmetrieachse & Dort kam $q_0$ aus der
    Verträglichkeit; der Verlauf ist trotzdem richtig \\
  Verwölbung nicht berechnet & Massiv, gemischt, mehrzellig, $I_T = 0$ oder Schubmittelpunkt
    unbestimmt \\
  Warnung: Offenes Profil, nicht in der Formelsammlung & $u_x = -\vartheta\int r\,ds + c$ ist
    Standardtheorie, aber nicht Teil der Formelsammlung \\
  Schub: Trägheitsmomente singulär & $I_yI_z - I_{yz}^2 \approx 0$ \\
  Hinweis: Hauptachsen nicht parallel, $I_{yz}$-Kopplung & Schiefe Biegung, siehe oben \\
  Torsion: nur rein dünnwandige Profile & Keine Prandtl-Lösung \\
  Torsion: mehrzelliges Profil & Statisch unbestimmt \\
  Torsion: $I_T = 0$ & Kein Torsionswiderstand \\
  Torsion aus Kräften: Schubmittelpunkt statisch unbestimmt & Das Moment ist nicht bestimmbar \\
  $\sigma_x$ nicht berechenbar (Nenner 0) & Fläche oder $I_yI_z - I_{yz}^2$ ist null \\
  Ungültige Eingabe & Der Ausdruck ließ sich nicht auswerten \\
  Hinweis: Profilbeiwert $\xi = 1$ -- so gewollt? & Torsion eines offenen Profils mit $\xi = 1$;
    die Formelsammlung nennt z.\,B.\ L 0,99, C/T 1,12, I 1,31 \\
  Hinweis: alle Elemente in der Standarddicke & Kommt einmal bei \taste{e}, wenn alle dünnwandigen
    Elemente die Standarddicke haben; erst wieder, wenn das Profil ganz gelöscht wurde \\
  Kein Querschnitt vorhanden & Es wurde noch nichts gezeichnet \\
  \taste{v} nur im Schubspannungsmodus & Erst Schub rechnen, dann Verlauf wählen \\
  Wirklich alles loeschen? & Rückfrage bei \taste{c}: \taste{Enter} löscht, \taste{Esc} bricht ab \\
  \bottomrule
\end{longtable}

\section{Was das Querschnittsskript nicht kann}

\begin{itemize}
  \item Kein Schub, keine Torsion und keine Verwölbung für mehrzellige Profile.
  \item Keine Prandtlsche Torsionsrechnung für massive Querschnitte.
  \item Keine Wölbkrafttorsion (die Verwölbung selbst schon, siehe oben), keine
        Querschnittsverformung, keine lokalen Spannungsspitzen, keine Kerbwirkung.
  \item Dünnwandige Kreise und Bögen gehen in die Querschnittswerte ein, nicht aber in die
        Schubrechnung.
  \item Numerisch addierte Elemente sind eine Rasternäherung (Warnung im Ergebnisfeld).
  \item Ein Material gibt es nicht; gerechnet wird rein geometrisch, $E$ kürzt sich heraus.
\end{itemize}

\section{Wie gerechnet wird}

Flächen, Schwerpunkte und Flächenmomente kommen aus geschlossenen Formeln je Grundform (Polygone
über die Gaußsche Trapezformel, Kreisausschnitte und -abschnitte über die entsprechenden
Integrale), zusammengesetzt mit dem Satz von Steiner. Dünnwandige Elemente werden als Mittellinie
mit Dicke behandelt -- genau wie in der Vorlesung.

Für den Schub wird aus den dünnwandigen Elementen ein Profilgraph gebaut: Die Elemente werden an
Kreuzungen und Berührpunkten geteilt, der Schubfluss startet an den freien Enden (dort ist er
null) und wird aufsummiert; an Verzweigungen gilt das Knotengleichgewicht. Eine geschlossene
Zelle wird aufgeschnitten, und der konstante Anteil $q_0$ kommt aus der Symmetrie. Massive
Querschnitte werden über waagerechte bzw.\ senkrechte Schnitte mit $\tau = QS/(Ib)$ ausgewertet.

%=====================================================================
\part{Spannungskreis}

\section{Die Idee: Eingabe über Schnitte}

In \textbf{TM2} ist die Aufgabe fast nie \enquote{gegeben sind $\sigma_x$, $\sigma_y$,
$\tau_{xy}$}, sondern \enquote{in diesem Schnitt ist $\sigma$ bekannt, in jenem $\tau$, und der
eine Winkel ist gesucht}. Genau so gebt ihr hier ein: in einer Tabelle, deren \textbf{Spalten die
Schnitte} sind und deren Zeilen

\begin{enumerate}
  \item der Winkel $\alpha$ der Schnittnormalen gegen die $x$-Achse, \textbf{mathematisch
        positiv} gezählt (gegen den Uhrzeigersinn) -- das steht auch über der Tabelle,
  \item die Normalspannung $\sigma$ in diesem Schnitt,
  \item die Schubspannung $\tau$ in diesem Schnitt
\end{enumerate}

sind. \textbf{Was ihr nicht wisst, lasst ihr leer} -- das Programm rechnet es aus, sobald genug
Angaben da sind. Ihr könnt also genauso gut einen Schnitt eingeben, von dem ihr nur $\sigma$
kennt, oder einen, von dem ihr $\sigma$ und $\tau$ kennt, aber nicht den Winkel.

Eingabe und Ergebnis sind dieselbe Tabelle: \textbf{schwarz} steht, was ihr eingegeben habt,
\textcolor{green!50!black}{\textbf{grün}}, was daraus folgt. Ihr könnt jederzeit in jede Zelle
zurück und den Wert ändern -- es wird sofort neu gerechnet. Rechts neben dem letzten Schnitt
steht eine Spalte \emph{neu}: schreibt ihr dort etwas hinein, entsteht ein weiterer Schnitt.

Unter der Tabelle sagt eine Zeile, woran ihr seid: \emph{eindeutig bestimmt},
\emph{noch $n$ unabhängige Angabe(n) nötig} oder \emph{Eingaben widersprechen sich}. Solange der
Zustand nicht eindeutig ist, sind \taste{s} und \taste{k} gesperrt -- ihr bekommt also kein Bild,
das gar nicht bestimmt ist.

\section{Bedienung}

\begin{longtable}{p{0.22\linewidth}p{0.74\linewidth}}
  \toprule
  \textbf{Taste} & \textbf{Wirkung} \\
  \midrule \endhead
  Pfeiltasten & Tabelle: Zelle wählen (links/rechts der Schnitt, oben/unten $\alpha$, $\sigma$,
    $\tau$); Kreis: Marke bewegen (links/rechts der Schnitt, oben/unten $5^\circ$) \\
  Ziffer, \code{.}, \code{-} & fängt sofort eine neue Eingabe in der Zelle an \\
  \taste{Enter} & Zelle zum Ändern öffnen bzw.\ übernehmen; danach geht es
    $\alpha \rightarrow \sigma \rightarrow \tau \rightarrow$ nächster Schnitt \\
  \taste{Backspace} & beim Ändern ein Zeichen löschen, sonst die Zelle leeren (der Wert wird dann
    wieder berechnet) \\
  \taste{Tab} & auf einem berechneten Winkel mit zwei Lösungen: zur anderen Lösung wechseln \\
  \taste{b} & rechnen \\
  \taste{s} & Spannungsscheibe \\
  \taste{k} & Mohrscher Spannungskreis \\
  \taste{t} & zurück zur Tabelle \\
  \taste{Del} & alles löschen, mit Rückfrage (\taste{Enter} löscht, \taste{Esc} bricht ab) \\
  \bottomrule
\end{longtable}

In die Zellen dürft ihr auch rechnen: \code{10*2}, \code{sqrt(2)} und \code{1,5} gehen alle.

\section{Die beiden Bilder}

\textbf{Spannungsscheibe (\taste{s}):} Gezeichnet wird im gewohnten Koordinatensystem:
\textbf{$x$ zeigt nach unten, $y$ nach rechts} -- genau wie im Querschnittsskript. Ein Schnitt mit
$\alpha = 0$ liegt damit unten. Die Scheibe ist ein \textbf{Vieleck mit einer Kante je
Schnitt} -- also genau die Figur, die ihr in der Aufgabe seht. An jeder Kante hängt ein roter
Pfeil für die Normalspannung (nach außen bei Zug, auf die Kante zu bei Druck) und ein blauer für
die Schubspannung längs der Kante, dazu die Beschriftung mit $\alpha$, $\sigma$ und $\tau$
außerhalb der Kante. Die Scheibe hat \textbf{genau so viele Kanten, wie ihr Schnitte eingegeben
habt} -- ein Trapez aus vier Schnitten bleibt ein Trapez. Nur wenn eure Schnitte die Scheibe nicht
schließen (bei zwei Schnitten zum Beispiel), ergänzt das Programm den $x$- oder $y$-Schnitt, und
zwar nur die Seite, die zum Schließen gebraucht wird: Aus 0° und 90° wird so das übliche Rechteck,
aus 30° und 120° ein Dreieck mit dem $y$-Schnitt. Diese Kanten heißen \code{x-Schnitt} bzw.\
\code{y-Schnitt}, ihre Spannungen kommen aus dem Zustand. Längliche Scheiben werden verkleinert,
bis sie ins Bild passen. Gezeichnet wird so, wie ihr es von Hand macht: Zug zeigt immer von der Scheibe weg, auf beiden
gegenüberliegenden Seiten, und die Schubpfeile gegenüberliegender Kanten zeigen in
entgegengesetzte Richtungen -- bei positivem $\tau_{xy}$ laufen sie paarweise auf dieselben Ecken zu.

\textbf{Mohrscher Kreis (\taste{k}):} $\sigma$ nach rechts, $\tau$ nach oben, der Kreis um
$(\sigma_m, 0)$ mit Radius $R$. Die $\tau$-Achse steht dort, wo sie hingehört, nämlich bei
$\sigma = 0$ -- ihr seht also ein richtiges Koordinatensystem und könnt direkt sehen, ob der
Zustand ganz im Zug- oder Druckbereich liegt. Hauptspannungen, Mittelpunkt und Nullpunkt sind
markiert, jeder Schnitt ist als Punkt mit Radiusstrahl eingezeichnet. Liegt der Kreis sehr weit
vom Ursprung weg, hat der Kreis Vorrang; der Nullpunkt steht dann als gepunktete Linie am
Bildrand. Mit den \textbf{Pfeiltasten} wandert ihr über den Kreis: links und rechts springen von
Schnitt zu Schnitt, hoch und runter drehen in $5^\circ$-Schritten weiter. Unten stehen dabei
immer $\alpha$, $\sigma$ und $\tau$ der Marke -- so bekommt ihr auch den Zustand in einem Schnitt,
den ihr gar nicht eingegeben habt. Denkt an den Faktor~2: Eine Drehung des Schnitts um $\alpha$ entspricht im Kreis
$2\alpha$, und bei wachsendem $\alpha$ wandern die Punkte im Uhrzeigersinn.

\bildplatz{spannungskreis}{Tabelle, Spannungsscheibe und Mohrscher Kreis}

\section{Wann ist der Zustand bestimmt?}

Mit $\sigma_m = \frac{\sigma_x+\sigma_y}{2}$, $R_c = \frac{\sigma_x-\sigma_y}{2}$ und
$R_s = \tau_{xy}$ ist

\[
  \sigma(\alpha) = \sigma_m + R_c\cos 2\alpha + R_s\sin 2\alpha, \qquad
  \tau(\alpha) = -R_c\sin 2\alpha + R_s\cos 2\alpha .
\]

Beides ist in $(\sigma_m, R_c, R_s)$ \emph{linear}, solange der Winkel bekannt ist. Jede Angabe
ist also eine lineare Gleichung, und \textbf{drei unabhängige Angaben} bestimmen den Zustand.
Das Programm löst dieses System und zählt dabei den Rang -- \enquote{bestimmt} meldet es erst,
wenn der Rang wirklich~3 ist. Zwei parallele Schnitte oder dieselbe Angabe zweimal bringen
nichts Neues, und das merkt es auch. Habt ihr mehr eingegeben als nötig, prüft es die Angaben
gegeneinander.

\textbf{Schnitte ohne Winkel} zählen auch mit, wenn ihr $\sigma$ \emph{und} $\tau$ kennt: Dann
liegt der Punkt $(\sigma, \tau)$ auf dem Kreis,

\[
  (\sigma - \sigma_m)^2 + \tau^2 = R_c^2 + R_s^2 ,
\]

und das ist genau eine Angabe (der Winkel ist ja die zweite Unbekannte). Typischer Fall ist
\emph{Klausur Aufgabe 5}: Hauptrichtung $x$ (eingeben als $\alpha = 0$, $\tau = 0$), dazu
$\sigma_y = 10$ ($\alpha = 90^\circ$) und der Schnitt $\sigma_\xi = 35$, $\tau_{\xi\eta} = 20$ ohne
Winkel. Das Programm liefert $\sigma_I = 51$, $\sigma_{II} = 10$ und den Winkel des Schnitts,
$-38{,}7^\circ$ (in der Tabelle als $141{,}3^\circ$, das ist dasselbe).

Weil die Gleichung quadratisch ist, kann es \textbf{zwei} passende Zustände geben -- etwa wenn ihr
$\sigma_x$, $\sigma_y$ und einen Punkt ohne Winkel kennt: Dann ist das Vorzeichen von $\tau_{xy}$
offen. Das Programm sagt dann \emph{zwei Zustände möglich}, zeigt beide an und lässt
\taste{s} und \taste{k} gesperrt, bis eine weitere Angabe entscheidet. Ein Schnitt ohne Winkel mit
nur $\sigma$ (oder nur $\tau$) ist zwar keine Gleichung, schließt aber einen Zustand aus, auf dessen
Kreis der Wert nicht passt.

\textbf{Gesuchte Winkel} folgen danach in geschlossener Form -- dafür braucht es kein CAS:

\begin{longtable}{p{0.24\linewidth}p{0.50\linewidth}p{0.18\linewidth}}
  \toprule
  \textbf{gegeben} & \textbf{Lösung} & \textbf{Anzahl} \\
  \midrule \endhead
  $\sigma$ und $\tau$ & $2\alpha = \varphi + \operatorname{atan2}(-\tau,\ \sigma-\sigma_m)$ &
    eindeutig \\
  nur $\sigma$ & $2\alpha = \varphi \pm \arccos\frac{\sigma-\sigma_m}{R}$ & zwei \\
  nur $\tau$ & $2\alpha = \varphi + \arcsin\frac{-\tau}{R}$ bzw.\ $\varphi + \pi - \arcsin\frac{-\tau}{R}$ & zwei \\
  \bottomrule
\end{longtable}

mit $\varphi = \operatorname{atan2}(R_s, R_c)$. Bei zwei Lösungen steht die zweite klein hinter
dem Winkel in der Tabelle. Mit \taste{Tab} auf diesem Feld wechselt ihr zur anderen Lösung; die
berechnete Spannung des Schnitts, die Scheibe und der Kreis ziehen mit. Das braucht ihr, wenn die
Aufgabe über eine Zusatzangabe (z.\,B.\ die Richtung der Schubspannung) festlegt, welche der beiden
Lösungen gemeint ist. Liegt euer Wert gar nicht auf dem Kreis (also
$|\sigma - \sigma_m| > R$), sagt euch das Programm das für diesen Schnitt.

\section{Grenzen und Tricks}

\begin{itemize}
  \item Schnitte ohne Winkel legen nie die \emph{Richtung} fest, nur Mittelpunkt und Radius --
        mindestens eine Angabe mit Winkel braucht ihr immer.
  \item Eine Hauptrichtung gebt ihr als Schnitt mit diesem Winkel und $\tau = 0$ ein. Ob dort
        $\sigma_I$ oder $\sigma_{II}$ wirkt, rechnet das Programm aus; eintragen lässt sich das nicht.
  \item Wenn ihr nur die Winkel \emph{zwischen} den Schnitten kennt: Setzt den ersten Schnitt auf
        $0^\circ$ und die anderen relativ dazu. Alle Ergebnisse sind dann auf diesen Schnitt
        bezogen, und die Hauptspannungen stimmen sowieso.
  \item Einen Schnitt werdet ihr los, indem ihr seine drei Zellen leert; \taste{Del} fängt
        komplett neu an.
  \item Nur ebener Spannungszustand, kein räumlicher und kein Verformungskreis.
\end{itemize}

Das ältere Modul \code{Mohr\_Spann.lua} gibt es weiterhin; dort gebt ihr fest vorgegebene
Kombinationen ein ($\sigma_x$, $\sigma_y$, $\tau_{xy}$, $\alpha$ / zwei Punkte / Punkt plus
Winkel zu $\tau_{max}$). Für alles, was nach Schnitten gefragt ist, ist das neue Modul schneller.


%=====================================================================
\appendix
\part{Anhang}

\section{Tastenkürzel Tragwerksskript}

\begin{longtable}{p{0.14\linewidth}p{0.82\linewidth}}
  \toprule
  \textbf{Taste} & \textbf{Funktion} \\
  \midrule \endhead
  \taste{b} & Tragwerk berechnen \\
  \taste{s} & Systemansicht \\
  \taste{n}, \taste{q}, \taste{m} & Normalkraft, Querkraft, Moment \\
  \taste{w}, \taste{u} & Biegelinie, Längsverformung \\
  \taste{e}, \taste{t} & Explosionsansicht, Beschriftungen dort \\
  \taste{v}, \taste{p} & Verschiebungszustand (PvV), Polplan \\
  \taste{l} & globales Gleichungssystem \\
  \taste{o} & Obermenü (Einstellungen und Exporte) \\
  \taste{h} & System einpassen \\
  \taste{f} & Stabrichtung umdrehen \\
  \taste{c}, \taste{Clear} & alles löschen (Rückfrage: \taste{Enter} löscht, \taste{Esc} bricht ab) \\
  \taste{+}, \taste{-} & Zoom \\
  \taste{1}\dots\taste{9} & KGV-Einheitszustand wählen (bei Auto-KGV) \\
  \taste{Enter} & Menü öffnen, Wert umschalten, Eingabe übernehmen \\
  \taste{Esc} & Box, Eingabe, Menü oder Modus schließen \\
  \taste{Backspace} & Eingabe oder Objekt löschen \\
  \bottomrule
\end{longtable}

\section{Tastenkürzel Querschnittsskript}

\begin{longtable}{p{0.14\linewidth}p{0.82\linewidth}}
  \toprule
  \textbf{Taste} & \textbf{Funktion} \\
  \midrule \endhead
  \taste{m}, \taste{d} & Massiv- bzw.\ Dünnwandmenü \\
  \taste{b} & Berechnungsmenü \\
  \taste{e}, \taste{t} & Querschnittswerte, Einzelwerte-Tabelle \\
  \taste{v} & Schubverlauf wählen (danach Ziffern \taste{0}\dots\taste{9}) \\
  \taste{k} & Kernflächenmodus weiterschalten \\
  \taste{o} & Optionen (Raster, Dicke, KOS, Einheiten, $\xi$, FTM-Bezug) \\
  \taste{h} & Geometrie einpassen \\
  \taste{c}, \taste{Clear} & alles löschen (Rückfrage: \taste{Enter} löscht, \taste{Esc} bricht ab) \\
  \taste{+}, \taste{-} & Zoom \\
  \taste{Enter} & Menüpunkt wählen, Element bearbeiten, Eingabe übernehmen \\
  \taste{Esc} & Box, Eingabe oder Ansicht schließen \\
  \taste{Backspace}, \taste{Del} & letztes Element löschen; in Eingabemasken ein Zeichen \\
  \bottomrule
\end{longtable}

\section{Tastenkürzel Spannungskreis}

\begin{longtable}{p{0.14\linewidth}p{0.82\linewidth}}
  \toprule
  \textbf{Taste} & \textbf{Funktion} \\
  \midrule \endhead
  Pfeiltasten & Zelle wählen; im Kreis die Marke bewegen \\
  Ziffer, \code{.}, \code{-} & Eingabe in der Zelle beginnen \\
  \taste{b} & rechnen \\
  \taste{s}, \taste{k} & Spannungsscheibe, Mohrscher Kreis (nur bei eindeutigem Zustand) \\
  \taste{t}, \taste{e} & zurück zur Tabelle \\
  \taste{Del}, \taste{Clear} & alles löschen (Rückfrage: \taste{Enter} löscht, \taste{Esc} bricht ab) \\
  \taste{Enter} & Zelle öffnen bzw.\ übernehmen \\
  \taste{Backspace} & Zeichen löschen bzw.\ Zelle leeren \\
  \taste{Tab} & andere Winkellösung (bei zwei Lösungen) \\
  \taste{Esc} & Eingabe abbrechen, sonst zurück zur Tabelle \\
  \bottomrule
\end{longtable}


\section{Häufige Stolpersteine}

\begin{itemize}
  \item \textbf{\code{t} und \code{T} gehen nicht als Symbol}, das ist die Integrationsvariable.
        Nehmt \code{dt} oder \code{t1}.
  \item \textbf{Griechische Buchstaben} werden nicht als Symbol erkannt; \code{a} statt $\alpha$.
  \item \textbf{Zahl und Symbol nicht mischen:} Wenn $EA$ symbolisch ist, sollten alle Stäbe ein
        symbolisches $EA$ haben; sonst kommen Terme wie \code{3*l\^{}2*ea + 125000000} heraus.
  \item \textbf{Pendelstäbe brauchen Gelenke} an beiden Enden, sonst sind es Biegestäbe mit dem
        voreingestellten $EI$.
  \item \textbf{Lager, an denen nur Stäbe mit Endgelenk hängen}, brauchen ein Knotengelenk; sonst
        ist die Knotenverdrehung unbestimmt und das System gilt als kinematisch.
  \item \textbf{Gleichmäßige Erwärmung:} $T_{oben} = T_{unten} = \Delta T$ eingeben.
  \item \textbf{Ein Bereich, ein Stab:} Überall, wo ihr von Hand eine neue Biegelinie ansetzen
        würdet (Lager, Gelenk, Lastsprung), braucht das Programm einen eigenen Stab.
  \item \textbf{Nach dem Export} lohnt ein Blick in \code{randinfo}: Dort steht, welche Funktionen
        Verschiebungen sind und welche $EI\,w$ bzw.\ $EA\,u$.
\end{itemize}


\end{document}
